\section{Introduction}

Kac's famous question --- can one hear the shape of a drum? --- remains a
useful way to frame inverse spectral problems, provided one resists the urge to
promise more than the data support \cite{Kac1966}.  In structural acoustics and
vibroelasticity, the practical issue is usually local rather than global: given
a finite set of measured resonances, can nearby changes in geometry and
stiffness be separated reliably?  That is a practical-identifiability question
in the usual inverse-problem sense \cite{Raue2009,Wieland2021}.  For that
purpose, the singular values of a scaled frequency--parameter Jacobian are the
natural diagnostic.

The Browntone companion papers established the ingredients needed for this
question.  Paper~1 introduced the canonical fluid-filled oblate shell used as
an abdominal benchmark \cite{BrowntoneP1}.  Paper~7 showed that an
equivalent-sphere reduction can be very effective for a forward problem in
which one seeks a dominant stiffness indicator rather than full geometric
recovery \cite{BrowntoneP7}.  Paper~8 then mapped local practical
identifiability across oblate and prolate shell families \cite{BrowntoneP8}.
The present short communication focuses on the corresponding asymmetry result:
\emph{oblateness lifts a weak spectral inverse direction, whereas prolateness
does not provide a comparable improvement}.

That distinction matters because ``symmetry breaking helps'' is false in the
form in which it is usually repeated.  The five-mode flexural map considered
here does not reward every departure from spherical geometry.  Instead, the
improvement is tied to how the perturbed geometry redistributes curvature across
the modal strain field.  Oblate shells separate the geometric sensitivity
directions enough to improve conditioning; prolate shells, in the same
framework, remain less helpful.

Using the canonical fluid-loaded shell and the flexural modes $n=2,\ldots,6$
derived from thin-shell theory and fluid--structure coupling
\cite{Lamb1882,JungerFeit1986}, we show that the inverse map moves from an
equivalent-sphere singular limit to a near-spherical Ritz floor and then to a
useful oblate operating point, while the prolate branch remains less well
conditioned.  The lowest flexural resonance in that map is
$f_2=\SI{3.95}{\hertz}$ at the canonical oblate point \cite{BrowntoneP1}, so
the discussion concerns the low-frequency flexural spectrum only; the much
higher breathing branch is outside the present scope.
